% id: 164-327
% book: 개념원리 미적분Ⅱ · 19 접선의 방정식
% section: 연습문제 STEP 2
% figure: none
% answer: ③
% answer_source: 답지
% variant_level: 0 (원본 전사)
% 이 파일은 build.mjs 가 items.json 에서 만든다. 고칠 때는 items.json 을 고친다.
\smash{\raisebox{1.5mm}{\dmkichul{개념원리 미적분Ⅱ 164쪽 327번 · 평가원 기출}}}
실수 $t$ $(0<t<\pi)$에 대하여 곡선 $y=\sin x$ 위의 점~$\pt{P}(t{,}\mkern3mu\mathopen{}\sin t)$에서의 접선과 점~$\pt{P}$를 지나고 기울기가 $-1$인 직선이 이루는 예각의 크기를 $\theta$라 할 때, $\displaystyle\lim_{t\to\pi-}\dfrac{\tan\theta}{(\pi-t)^2}$의 값은?
\begin{choices32}{\rule[-3ex]{0pt}{8ex}$\dfrac{1}{16}$}{\rule[-3ex]{0pt}{8ex}$\dfrac{1}{8}$}{\rule[-3ex]{0pt}{8ex}$\dfrac{1}{4}$}{\rule[-3ex]{0pt}{8ex}$\dfrac{1}{2}$}{\rule[-3ex]{0pt}{8ex}$1$}\end{choices32}
